\section{Related Work}
\label{sec:related}

\paragraph{Frequency-domain detectors.}
The frequency-domain approach to fake-image detection predates the current wave of diffusion models. Durall~\etal~\cite{durall2020watch} showed that GAN upsampling layers leave a characteristic deviation from the $1/f^{\alpha}$ power law of natural images and used the azimuthally averaged power spectrum as a near-linearly-separable feature. Zhang~\etal~\cite{zhang2019detecting} proposed a spectrum-based classifier targeting GAN checkerboard artifacts, and Frank~\etal~\cite{frank2020leveraging} trained CNNs directly on DCT coefficients. The shared finding is that discrete upsampling operators leak energy into specific frequency bands in ways that generalise across generators sharing the same architectural primitives.

Diffusion models are different: they decode through a learned VAE rather than a stack of transposed convolutions. Corvi~\etal~\cite{corvi2023intriguing} and Ricker~\etal~\cite{ricker2024aeroblade} report that frequency-based signals partially transfer to the diffusion setting, though neither paper runs the controlled 1D-vs-2D comparison we pursue here. Bammey~\cite{synthbuster2023} and Karageorgiou~\etal~\cite{spai2025} push frequency-only detection to current diffusion benchmarks with competitive results; SPAI~\cite{spai2025} in particular uses masked spectral learning on real images only and outperforms several supervised baselines.

\paragraph{Semantic and learned-feature detectors.}
A second family of detectors uses pretrained visual encoders --- CLIP~\cite{radford2021clip}, DINOv2~\cite{oquab2023dinov2} --- as frozen feature extractors and trains a shallow classifier on top~\cite{ojha2023univfd,cozzolino2024raising}. These methods tend to generalise well across generators but are computationally heavier and harder to interpret than frequency-only approaches. Wang~\etal~\cite{wang2020cnn} is the canonical pixel-space CNN detector; our work is the frequency-domain analogue of that paper, applied to the diffusion era.

\paragraph{Robustness and dataset artifacts.}
Several recent papers show that detector performance on standard benchmarks can be confounded by post-processing. Grommelt~\etal~\cite{grommelt2024} demonstrate that GenImage~\cite{genimage} has a systematic JPEG quality difference between its real and fake classes, meaning a naive detector can score above 0.9 AUC while essentially classifying compression quality rather than generative origin. Any frequency-domain detector is especially vulnerable to this, because JPEG quantisation itself distorts the spectrum. We flag this explicitly when discussing our GenImage results (Section~\ref{sec:genimage}) and treat Defactify~\cite{defactify} as a less-confounded complement.

\paragraph{Invisible watermarking and SynthID.}
An orthogonal approach to provenance detection is to embed a watermark at generation time. Google's SynthID~\cite{synthid2023} applies an imperceptible frequency-domain perturbation to images produced by Gemini and related systems; the watermark persists through moderate JPEG compression and resizing and can be decoded without the original image. Unlike most prior watermarking schemes~\cite{cox2007watermarking}, SynthID is designed to be robust to cropping and colour adjustments. The standard detection pipeline requires the watermarking key. Our work is complementary: we ask whether the \emph{presence} of the watermark is observable as a spectral anomaly, without access to the key (Section~\ref{sec:synthid}).

\paragraph{Adversarial attacks on detectors.}
PGD~\cite{madry2018pgd} is the standard first-order method for evaluating classifier robustness. Most prior attacks on AI-image detectors operate in pixel space~\cite{carlini2020evading}; because our detector's input \emph{is} the Fourier transform of the image, a pixel-space attack and a frequency-space attack are related by a linear change of variables but differ in which perturbations look natural. We do not run adversarial experiments in this project but include this context because frequency-domain detectors are a natural target for such attacks and robustness analysis is a clear direction for future work.

\paragraph{Datasets.}
CIFAKE~\cite{cifake} pairs CIFAR-10 photographs with Stable Diffusion v1.4 images of the same classes and has become a standard small-scale benchmark. GenImage~\cite{genimage} covers eight generators (ADM, BigGAN, GLIDE, Midjourney, SD v1.4, SD v1.5, VQDM, Wukong) with ImageNet real photographs; we use four of the eight due to extraction and inode constraints described in Section~\ref{sec:genimage}. GenImage++~\cite{genimagepp} extends the protocol with harder, more recent generators and is intended as a test-only holdout. Defactify~\cite{defactify} is distributed as HuggingFace parquet shards and contains images from several modern generators. FaceForensics++~\cite{ffpp} is the dominant face-manipulation benchmark and is largely orthogonal to the present work.
